\documentclass[5pt]{article}
\usepackage{multicol,multirow}
\usepackage{graphicx} % Required for inserting images
\usepackage[margin=0.75cm]{geometry}
\usepackage{xcolor}
\usepackage{amsmath,esint}
\usepackage{mathtools}
\usepackage{relsize}
\usepackage{mathtools}
\usepackage{nccmath}
\usepackage[inline]{enumitem}
\usepackage{algpseudocode}
\usepackage{physics}
\usepackage[frak=esstix]{mathalpha}
\usepackage{array}

\usepackage{empheq}
\usepackage{amsfonts}

\definecolor{LightGray}{gray}{0.9}

\usepackage{minted}

\newenvironment{amatrix}[1]{%
  \left[\begin{array}{@{}*{#1}{c}|c@{}}
}{%
  \end{array}\right]
}


\makeatletter
\let\oldabs\abs
\def\abs{\@ifstar{\oldabs}{\oldabs*}}

\newtheorem{postulates}{Postulates}
\newtheorem{theorem}{Theorem}
\newtheorem{properties}{Properties}

\newcolumntype{L}{>{\footnotesize}l}
\newcolumntype{C}{>{\ttfamily\footnotesize}l}

\begin{document}

\begin{center}
     \Large{\textbf{Network Principles}}\\
     \footnotesize{Class: COSC 2670}\hfill\footnotesize{Maximilien Notz}
     \noindent\rule{20.2cm}{0.4pt}
\end{center}


\begin{multicols}{2}
\setcounter{secnumdepth}{0}

\subsection{Definitions}
\begin{tabular}{LL}
  IPv4                & The most used internet protocol, 32-bit binary value. \\
  IPv6                & The newest internet protocol, 128 bits IP adresse. \\
  IP Address          & A unique number identifier for a device on a \\
                      & network. Represented as 4 octets separated by dots. \\
                      & Example: \texttt{8.10.255.1} \\
  MAC Address         & \\
                      & a network. To make it readable the IP is represented. \\
  Domain Name         & Human readable string of character that map to a\\
                      & IP address. \\
  DNS                 & The DNS (Domain Name System) translates domain\\
                      & names to IP addresses. \\
  ISP                 & Internet Service Provider. Company that provides\\
                      & internet access. \\
  External IP Address & Assigned to a Router by the ISP. This address is\\
                      & unique across internet. \\
  Internal IP Address & Assigned to a device by the Router. Unique within\\
                      & the network. \\
  NAT                 & Network Address Translation, Allows 2 device to\\
                      & share the same external IP address. \\
  RJ11                & Telephone cable, stand for Registered Jack-11.\\
  RJ45                & Ethernet cable, stand for Registered Jack-45.\\
  Switch              & Switchs connect multiple computers by RJ45. \\
  LAN                 & Local Area Network, . \\
  Router              & Connects a LAN to the internet. More generally,\\
                      & connects different networks together. \\
  network diagram     & A visual representation of a network's structure and\\
                      & components. \\
\end{tabular}


\subsubsection{Network Address Translation (NAT)}
  The router takes the requests of the devices whithin the network and send it to the internet. 
  When the answer comes back the Router redistribute the packet to the devices.
  \textbf{Security Benefit:} Devices with internal IPs are hidden from public internet.

\subsection{Network Diagrams}
\subsubsection{Golden rules}
\begin{itemize}
  \item Never set up a network installation without first creating a network diagram.
  \item Never modify your network without first consulting and updating your network diagram.
  \item You should always have your network diagram at hand.
\end{itemize}

\subsection{Internet Protocol (IP)}
Internet Protocol, responsible for exchanging data between devices on the network.
\subsubsection{IP Packets}
The data exchanged between machines is split into packets of generally 1500 bytes. To be sent over the network sequentially. 
The data is rebuilt at the receiving end machine. 
\textbf{Packet Header contains:} A source IP Address, A destination IP Address, The data being sent, the size of the packet and other information depending on the protocol.


\subsection{Routing}



\subsection{Linux Networking Commands}
\begin{tabular}{C L}
  ifconfig    & Displays IP address.\\
  ping        & Returns packet loss.\\
  traceroute  & Traces the path packets take to a destination.\\
  nslookup    & Returns IP address information for a domain name.\\
  arp -a      & Displays the ARP table, which maps IP addresses to\\
              & MAC addresses.\\
  arp myIPa   & Displays the ARP table for the specific IP address .\\
  curl        & Transfers data from or to a server, using various\\
              & protocols.\\
  curl ifconfig.me    & Display machine's public IPv4 address.\\
\end{tabular}

\end{multicols}
\end{document}
